% =============================================================================
% LaTeX preamble for nox-mem paper (arXiv-safe, xelatex)
%
% Monospace: default Latin Modern Mono (from lmodern) -- bundled with every
% TeX Live, including arXiv's AutoTeX. We deliberately do NOT \setmonofont a
% system font: Courier New / DejaVu Sans Mono are absent on arXiv's Linux
% TeX Live and would break the xelatex build (fontspec "font not found").
%
% Trade-off: a handful of rare Unicode math glyphs (e.g. inside code blocks)
% may fall back to a box; these are compiler warnings, not errors, and prose
% math renders correctly via unicode-math.
% =============================================================================

\usepackage{fontspec}
\usepackage{amssymb}
\usepackage{amsmath}

% ---------------------------------------------------------------------------
% Compact wide tables: shrink all longtables + tighten column padding so
% multi-column comparison tables (e.g. the 9-column "Six Gaps") don't squeeze
% a prose column down to one word per line. All arXiv-safe (etoolbox in TeX Live).
% ---------------------------------------------------------------------------
\usepackage{etoolbox}
\AtBeginEnvironment{longtable}{\footnotesize}
\setlength{\tabcolsep}{4pt}

% ---------------------------------------------------------------------------
% A CAUSA de todo glifo perdido: `fontspec` carregado e nunca instruido.
%
% O template do pandoc faz \usepackage{fontspec} sob XeTeX mas NAO chama
% \setmainfont, e sem isso o fontspec mantem a Latin Modern *Type1/EC*, de 256
% slots. Medido em 2026-09-11 neste manuscrito: 63 "Missing character" em
% ec-lmr10, ec-lmr8, ec-lmbx8, ec-lmtt8, ec-lmtt10 — os glifos alpha, seta,
% en dash, em dash, reticencias, menos, conjunto vazio, intersecao, aproximado
% e maior-ou-igual. Cada um **desaparece do PDF sem erro**: o build sai 0 e o
% caractere nao esta la.
%
% A Latin Modern OpenType traz todos e vem na TeX Live. Carregada pelo NOME DO
% ARQUIVO porque por nome de familia o fontspec nao a encontra nesta
% instalacao ("The font ... cannot be found").
%
% Isto substitui quatro remendos que eu tinha empilhado um a um, cada um a
% revelar o glifo seguinte: era UM defeito, nao cinco.
% ---------------------------------------------------------------------------
\setmainfont{lmroman10-regular.otf}[
  BoldFont       = lmroman10-bold.otf,
  ItalicFont     = lmroman10-italic.otf,
  BoldItalicFont = lmroman10-bolditalic.otf]
\setsansfont{lmsans10-regular.otf}[
  BoldFont   = lmsans10-bold.otf,
  ItalicFont = lmsans10-oblique.otf]
\setmonofont{lmmono10-regular.otf}[
  BoldFont   = lmmonolt10-bold.otf,
  ItalicFont = lmmono10-italic.otf,
  Scale      = MatchLowercase]

% ---------------------------------------------------------------------------
% Um glifo que nem a Latin Modern OpenType tem em modo texto: o alpha de
% "Mem-alpha", NOME de um sistema citado — soletra-lo seria errar o nome de
% quem se cita. Vem da Latin Modern Math, com `\char` porque o caractere e'
% ativo e reentraria.
%
% Os outros cinco (conjunto vazio, intersecao, uniao n-aria, aproximadamente,
% maior-ou-igual) foram reescritos em ASCII no fonte, e nao por preferencia: o
% construto `{\glifomat\char"..}` QUEBRA dentro de `\texttt{}` no documento
% completo ("Missing $ inserted"), embora compile isolado. Tres deles estavam
% exatamente la, num predicado de teoria de conjuntos num code span.
% ---------------------------------------------------------------------------
\newfontfamily\glifomat{latinmodern-math.otf}
\catcode`\α=\active \defα{{\glifomat\char"03B1}}
